
\begin{abstract}
We continue an audit-driven program in which finite-resolution scanning and readout are treated as primary mathematical structure, rather than postulates external to an ``ontic'' dynamics. Building on an earlier closed pipeline from noncommutative scanning to cusp coefficient spectra, Hecke dynamics, and Langlands semantics, we introduce an additional layer ``above the cusp'': a period--motive interface that explains why stable constants repeatedly appear as $\pi$, $\log$, and zeta-values.

On a controlled class of protocols we define a category $\ScanAlg$ whose objects are Kronecker-type scans on tori equipped with algebraic (rational) readout kernels and auditable regularization rules. We construct a \emph{holographic scanning functor} $\HSP:\ScanAlg\to\PerDatum$ into a category of period data and prove a closed realization theorem: for rationally independent scan slopes, the long-time Birkhoff readout equals the Kontsevich--Zagier period associated with $\HSP(\mathcal{P})$. Finite-resource implementations admit an explicit auditable error decomposition into (i) a sampling discrepancy term controlled by diophantine properties of the scan and (ii) a truncation/regularization term with provable bounds.

We provide reproducible pure-Python experiments realizing $\log 2$, $\pi$, $\zeta(2)$ and $\zeta(3)$ via one-, two-, and three-dimensional scans using truncated geometric kernels. As a toy model for ``constant selection'' under a bounded description budget, we reproduce a low-complexity search in which $4\pi^3+\pi^2+\pi$ emerges as the unique best approximation to $\alpha^{-1}$ among nonnegative integer combinations of $\pi,\pi^2,\pi^3$ with bounded coefficient sum. Finally, we formulate a falsifiable selection principle: observable constants correspond to period data that minimize a stability--complexity functional under finite resource constraints, motivating a minimal operative notion of \emph{wish} as a protocol-invariant, stably realizable period structure.
\end{abstract}

\begin{sloppypar}
\textbf{Keywords}: holographic scanning principle; Kronecker scan; equidistribution; discrepancy; Koksma--Hlawka inequality; Kontsevich--Zagier periods; motives; multiple zeta values; regularization; error budget; functoriality; selection principle.
\end{sloppypar}

\paragraph{Conventions.}
Unless otherwise stated, $\log$ denotes the natural logarithm. ``mod $1$'' refers to reduction in $\RR/\ZZ$. The scan time $t\in\ZZ_{\ge 0}$ denotes an iteration count of a protocol-defined update (Layer~1). We identify the $d$-torus with $\TT^d=(\RR/\ZZ)^d$ and use the fundamental domain $[0,1)^d$ when writing integrals.
When convenient, we freely replace $[0,1)^d$ by $[0,1]^d$ in integrals and discrepancy definitions, since the boundary has Lebesgue measure zero.

\tableofcontents
\newpage


